\documentclass[11pt,a4paper]{article}

% ---------------------------------------------------------------- packages
\usepackage[a4paper,margin=2.4cm,top=3.6cm,bottom=3cm]{geometry}
\usepackage[utf8]{inputenc}
\usepackage[T1]{fontenc}
\usepackage{lmodern}
\usepackage{amsmath,amssymb,amsthm}
\usepackage{graphicx}
\usepackage{xcolor}
\usepackage{enumitem}
\usepackage{fancyhdr}
\usepackage{titlesec}
\usepackage[most]{tcolorbox}
\usepackage{listings}
\usepackage{needspace}
\usepackage[colorlinks=true,linkcolor=navy,citecolor=navy,urlcolor=navy]{hyperref}
\usepackage{booktabs}

% ---------------------------------------------------------------- UCLouvain colours (matching the course slides)
\definecolor{navy}{HTML}{00214D}
\definecolor{ucblue}{HTML}{5DB3E6}
\definecolor{ucgrey}{HTML}{8591A0}
\definecolor{ucblock}{HTML}{F2F7FB}
\definecolor{ucorange}{HTML}{E07B00}
\definecolor{ucred}{HTML}{C1272D}

% ---------------------------------------------------------------- section styling
\titleformat{\section}
  {\color{navy}\Large\bfseries}{\thesection}{0.6em}{}
\titleformat{\subsection}
  {\color{navy}\large\bfseries}{\thesubsection}{0.6em}{}
\titlespacing*{\section}{0pt}{1.4em}{0.6em}

% ---------------------------------------------------------------- "Part" headers
\newcommand{\homeworkpart}[2]{%
  \vspace{0.6em}
  \begin{tcolorbox}[colback=navy,colframe=navy,arc=0pt,boxrule=0pt,
    left=10pt,right=10pt,top=7pt,bottom=7pt]
    {\color{white}\bfseries\large Part #1 --- #2}
  \end{tcolorbox}
  \vspace{0.4em}
}

% ---------------------------------------------------------------- task / deliverable boxes
\newcounter{taskcounter}
\newtcolorbox{task}{%
  enhanced, breakable, colback=ucblock, colframe=ucblue,
  boxrule=0.8pt, arc=2pt, left=8pt, right=8pt, top=6pt, bottom=6pt,
  before skip=8pt, after skip=8pt,
  fonttitle=\bfseries\color{navy},
  title={Task \stepcounter{taskcounter}\thetaskcounter}
}
\newtcolorbox{note}{%
  enhanced, breakable, colback=white, colframe=ucgrey,
  boxrule=0.6pt, arc=1pt, left=8pt, right=8pt, top=5pt, bottom=5pt,
  before skip=6pt, after skip=6pt,
  fonttitle=\bfseries\color{ucgrey}, title={Note}
}

% ---------------------------------------------------------------- code listings
\lstdefinestyle{code}{
  basicstyle=\ttfamily\small,
  keywordstyle=\bfseries\color{navy},
  commentstyle=\itshape\color{ucgrey},
  stringstyle=\color{ucorange},
  numbers=left, numberstyle=\tiny\color{ucgrey}, numbersep=8pt,
  breaklines=true, breakatwhitespace=false,
  frame=single, rulecolor=\color{ucblue!60}, framesep=6pt,
  backgroundcolor=\color{ucblock},
  showstringspaces=false, tabsize=2,
  xleftmargin=1em, xrightmargin=1em,
}
\lstset{style=code}

% ---------------------------------------------------------------- header / footer (mirrors the slide deck footer)
\pagestyle{fancy}
\fancyhf{}
\renewcommand{\headrulewidth}{0pt}
\renewcommand{\footrulewidth}{0.6pt}
\fancyfoot[L]{\small\color{navy}\bfseries LMECA2170 --- Computational Geometry}
\fancyfoot[R]{\small\color{ucgrey}\thepage\ /\ \pageref{LastPage}}
\usepackage{lastpage}

% ---------------------------------------------------------------- title block
\newcommand{\ucletterhead}{%
  \includegraphics[height=1.0cm]{figures/uclouvain.png}\\[8pt]
  {\color{navy}\rule{\linewidth}{1.6pt}}\\[10pt]
  {\color{ucblue}\rule{\linewidth}{0.8pt}}\\[14pt]
}

\begin{document}
\thispagestyle{fancy}

\ucletterhead

\begin{center}
  {\color{navy}\Huge\bfseries Computational Geometry}\\[6pt]
  {\color{navy!80}\Large Homework --- Delaunay Triangulation, Alpha Shapes \& PFEM}\\[10pt]
  {\large LMECA2170 \quad$\cdot$\quad Academic year 2026--2027}\\[8pt]
  {\color{ucgrey}Miguel Delecourt \(\cdot\) Vincent Legat \(\cdot\) Kerem Okyay \(\cdot\)
  Illy Perl \(\cdot\) Alexandre Pirot \(\cdot\)
  Jean-Fran\c{c}ois Remacle \(\cdot\) Felix Ruyfellaere \(\cdot\) Nathan Tihon}\\[2pt]
  {\color{ucgrey}Universit\'e catholique de Louvain}
\end{center}

\vspace{0.6em}
{\color{ucblue}\rule{\linewidth}{0.8pt}}
\vspace{0.8em}

\noindent
This homework has three parts, each building directly on the previous one. In
\textbf{Part~1} you implement, entirely from scratch, an incremental Delaunay
triangulation on top of a half-edge data structure, and filter it into an
\emph{alpha shape}. In \textbf{Part~2} you improve that alpha shape by
inserting Steiner points with \textbf{Chew's algorithm}, reusing the exact
same insertion routine you wrote in Part~1. In \textbf{Part~3} you assemble
everything into a \textbf{PFEM} (Particle Finite Element Method) time-stepping
loop for free-surface flows, in which the alpha shape of Part~1 is recomputed
at \emph{every time step} to track the moving fluid domain.

\begin{note}
Read the whole statement once before starting to code: the three parts share
data structures and routines, and the way you design Part~1 will determine
how painless (or painful) Parts~2 and~3 turn out to be.
\end{note}

% ======================================================================
\homeworkpart{1}{From a Point Cloud to an Alpha Shape}
% ======================================================================

Let $S=\{p_1,\dots,p_n\}\subset\mathbb{R}^2$ be a cloud of $n$ points. The
goal of Part~1 is to compute the Delaunay triangulation $DT(S)$
incrementally, using a \textbf{half-edge} mesh data structure that you
implement yourself, and then to filter $DT(S)$ into an \textbf{alpha shape}:
a subset of the triangles that approximates the ``shape'' formed by the
points, rather than just their convex hull.

\subsection*{1.1 --- The half-edge data structure}

You are free to use any half-edge (or half-edge-like) layout, but we
recommend the classic \emph{triangle-based} scheme, since it is the
simplest one to keep synchronized while triangles are created and destroyed:

\begin{itemize}[leftmargin=1.4em]
  \item Triangles are stored as a flat array \lstinline{triangles[3*nt]}: the
    three vertex indices of triangle $t$ are
    \lstinline{triangles[3t]}, \lstinline{triangles[3t+1]}, \lstinline{triangles[3t+2]},
    listed \textbf{counter-clockwise}.
  \item Each triangle $t$ owns exactly 3 half-edges, indexed $3t,3t+1,3t+2$.
    Half-edge $3t+k$ goes from vertex \lstinline{triangles[3t+k]} to vertex
    \lstinline{triangles[3t+(k+1)%3]}.
  \item \lstinline{next(he)} needs no storage: it is simply
    \lstinline{3*(he/3) + (he+1)%3}.
  \item One array \lstinline{opposite[3*nt]} stores, for every half-edge, the
    index of its \textbf{twin} (the half-edge of the neighbouring triangle
    that runs along the same edge in the opposite direction), or $-1$ if the
    edge is on the boundary.
\end{itemize}

This representation is compact, cache-friendly, and --- crucially for this
homework --- \emph{local}: inserting a point only touches the half-edges of
the triangles being removed and created, never the whole mesh.

\begin{task}\label{tsk:halfedge}
Implement the half-edge structure above (or an equivalent one of your
choice) together with the elementary queries you will need: the two
triangles sharing an edge, the three neighbours of a triangle, and the
vertices of a triangle in CCW order.
\end{task}

\subsection*{1.2 --- Bootstrapping: three points at infinity}

Incremental insertion needs \emph{some} triangulation to start from. Resist
the temptation to fake this with 3 extra vertices placed ``far enough away''
at some large finite coordinate: how far is far enough depends on $S$, gets
this wrong silently (a triangulation that is simply \emph{wrong}, not a
crash), and re-introduces exactly the floating-point magnitude problems this
course spends so much time warning you about. The correct construction uses
3 genuine \textbf{points at infinity} $q_0,q_1,q_2$, and adapts
\lstinline{orient2d}/\lstinline{incircle} to handle them --- exactly,
with no tuning parameter at all.

Represent each $q_k$ not by coordinates but by a \textbf{direction}
$d_k\in\mathbb{R}^2$ (e.g.\ three unit vectors $120^\circ$ apart). The
starting triangulation is the single triangle $(q_0,q_1,q_2)$: since the
three directions span the whole circle, every finite point is, by
convention, inside it --- no test is ever needed to check this.

\paragraph{Extending orient2d.} Write $A,B,C$ for the (finite or infinite)
vertices of the predicate. If a vertex is infinite with direction $d$, its
homogeneous row is $(d_x,d_y,0)$ instead of $(x,y,1)$; expanding the usual
$3\times3$ determinant with such a row gives the following closed forms
(each derivable from the others by the cyclic symmetry
$\text{orient2d}(A,B,C)=\text{orient2d}(B,C,A)$):

\begin{center}
\begin{tabular}{@{}ll@{}}
\toprule
Case & Formula \\
\midrule
$A,B,C$ all finite & the usual determinant \\
$C=\infty$, direction $d$ & $\text{orient2d}(A,B,C)=\mathrm{cross}(B{-}A,\,d)$ \\
$B=\infty$, direction $d$ & $\text{orient2d}(A,B,C)=\mathrm{cross}(d,\,C{-}A)$ \\
$A=\infty$, direction $d$ & $\text{orient2d}(A,B,C)=\mathrm{cross}(C{-}B,\,d)$ \\
$B,C=\infty$, directions $d_B,d_C$ & $\text{orient2d}(A,B,C)=\mathrm{cross}(d_B,d_C)$ (constant, any finite $A$)\\
\bottomrule
\end{tabular}
\end{center}

where $\mathrm{cross}(u,v)=u_xv_y-u_yv_x$. No division, no large
coordinates, no loss of precision --- ever.

\paragraph{Extending incircle.} The point $D$ being inserted is always
finite; $A,B,C$ are the vertices of the candidate triangle (listed CCW).

\begin{itemize}[leftmargin=1.4em]
  \item \textbf{0 infinite vertices}: the usual \lstinline{incircle}
    determinant.
  \item \textbf{1 infinite vertex}, say $C=\infty$ in direction $d$: as $C$
    recedes to infinity the circumcircle of $A,B,C$ degenerates into the
    \emph{half-plane} bounded by line $AB$, on the side $d$ points to. Hence
    $$\text{incircle}(A,B,C,D) \ \text{has the same sign as}\ \text{orient2d}(A,B,D).$$
    (Cyclically rotate if $A$ or $B$ is the infinite one instead: e.g.\ if
    $A=\infty$, use $\text{orient2d}(B,C,D)$.)
  \item \textbf{2 infinite vertices}, say $B=\infty$ (direction $d_B$) and
    $C=\infty$ (direction $d_C$), with $A$ finite: the circumcircle
    degenerates into the infinite \emph{angular wedge}, based at $A$, swept
    counter-clockwise from $d_B$ to $d_C$. Hence
    $$\text{incircle}(A,B,C,D) > 0 \iff \text{orient2d}(A,A{+}d_B,D)>0 \ \text{ and }\ \text{orient2d}(A,D,A{+}d_C)>0.$$
  \item \textbf{3 infinite vertices}: only the very first (bootstrap)
    triangle; it is never tested, since $S$'s first point is inside it by
    construction.
\end{itemize}

\begin{note}
Every case above reduces to one or two calls to \lstinline{orient2d} ---
you never need to evaluate a circumcircle when an infinite vertex is
involved. This is not a coincidence: it is exactly why using true points at
infinity is both simpler \emph{and} more robust than a large finite
super-triangle.
\end{note}

\subsection*{1.3 --- Insertion order: sorting along a Hilbert curve}

Insert the points of $S$ one at a time --- but not in an arbitrary order.
Sort $S$ along a \textbf{Hilbert space-filling curve} first. Consecutive
points along the curve are guaranteed to be close together in the plane, so
each new point almost always lands close to where the \emph{previous} one
was inserted, which is exactly what makes the point-location walk below fast
in practice (amortized $O(1)$ instead of $O(\sqrt n)$ per insertion).

\begin{task}
Implement a Hilbert-curve sort of the input points (map each point to a
Hilbert index at a fixed recursion depth, then sort by that index).
\end{task}

\subsection*{1.4 --- Locating the point: a stochastic visibility walk}

For every point $p$ (in Hilbert order), first find the triangle $T_0$ that
contains it, starting the walk from the triangle in which the
\emph{previous} point was inserted:

\needspace{12\baselineskip}
\begin{lstlisting}[style=code,caption={Stochastic visibility walk},captionpos=b]
T <- last triangle used
loop:
  let (u,v,w) be the vertices of T, CCW
  test the 3 edges (u,v),(v,w),(w,u) of T in a *random* order
  for the first edge (a,b) such that orient2d(a,b,p) < 0:
      cross into the triangle across (a,b);  T <- that triangle
      continue the loop
  # no edge tested negative: p is inside T
  T0 <- T;  stop
\end{lstlisting}

The predicate \lstinline{orient2d(a,b,p)} returns (twice) the signed area of
$abp$: negative means $p$ lies to the right of the oriented edge
$a\!\to\!b$, i.e.\ outside the triangle on that side.

\begin{note}
Randomizing the order in which the 3 edges are tested at every step is not
cosmetic: on a triangulation that is not Delaunay (which can happen
transiently, or if you experiment with non-Delaunay meshes), a
\emph{deterministic} edge order can make the walk cycle forever. The
randomized version almost surely terminates. On a genuine Delaunay
triangulation the walk always terminates, deterministic or not --- but get
into the habit of randomizing, it costs nothing and removes an entire class
of bugs.
\end{note}

\subsection*{1.5 --- Building the Delaunay cavity}

Starting from $T_0$, grow the \textbf{cavity} of $p$ breadth-first: a
triangle belongs to the cavity if and only if $p$ lies strictly inside its
circumcircle (the \lstinline{incircle} predicate).

\needspace{12\baselineskip}
\begin{lstlisting}[style=code,caption={Cavity computation (BFS)},captionpos=b]
cavity <- {T0};  Q <- [T0]
while Q is not empty:
   T <- Q.pop_front()
   for each neighbour T' of T, across shared edge e:
       if T' not in cavity and incircle(T', p) > 0:
           cavity <- cavity + {T'}
           Q.push(T')
       else:
           boundary <- boundary + {e}
\end{lstlisting}

At the end, \lstinline{boundary} is the ordered set of edges bounding the
star-shaped hole left by removing every triangle of \lstinline{cavity}.

\subsection*{1.6 --- Retriangulating the cavity}

Delete every triangle of \lstinline{cavity} and create one new triangle
$(a,b,p)$ for every boundary edge $(a,b)$, connecting $p$ to the whole
boundary loop (a \emph{fan}). This is the step where the half-edge structure
must be updated \textbf{locally}:

\begin{itemize}[leftmargin=1.4em]
  \item for each new triangle $(a,b,p)$, set the twin of its edge $(a,b)$ to
    whatever triangle used to be on the other side of that boundary edge
    (unchanged, since that neighbour was \emph{not} in the cavity);
  \item set the twins of the two new internal edges ($b\!\to\!p$ and
    $p\!\to\!a$) to their counterparts in the two adjacent new fan triangles.
\end{itemize}

No other half-edge in the mesh needs to be touched.

\begin{task}\label{tsk:bowyerwatson}
Implement the full incremental Bowyer--Watson loop: super-triangle,
Hilbert-sorted insertion, visibility walk, cavity, retriangulation. When
every point of $S$ has been inserted, discard every triangle that still
touches one of $q_0,q_1,q_2$: what remains is exactly $DT(S)$.
\end{task}

\subsection*{1.7 --- Filtering into an alpha shape}

An \textbf{alpha shape} of parameter $\alpha>0$ is obtained by removing,
from $DT(S)$, every triangle whose circumradius exceeds $\alpha$. Loosely,
$\alpha$ acts as a probe radius: triangles that could only exist because
Delaunay triangulates the \emph{convex hull} of $S$ --- bridging concavities
or holes with large, thin triangles --- have a large circumradius and
disappear, while the triangles that genuinely reflect the local density of
$S$ survive.

\begin{task}
Compute the circumradius of every triangle of $DT(S)$ and remove those
exceeding a user-supplied $\alpha$. Test on a point cloud shaped like a
non-convex outline (e.g.\ a letter, an annulus, or two point clusters joined
by a thin bridge) and check that the result visually matches the cloud's
shape for a reasonable range of $\alpha$, and degenerates back towards the
convex hull as $\alpha\to\infty$.
\end{task}

\paragraph{A single global $\alpha$ is not always enough.} If $S$ has a
strongly non-uniform density --- a dense cluster next to a sparse one --- no
single constant $\alpha$ works everywhere: small enough to avoid bridging
the sparse region, it eats away real detail in the dense one; large enough
to preserve the dense region, it bridges the sparse one anyway. The fix is
to let $\alpha$ become a \textbf{size map} $\alpha(x,y)$, a scalar field over
the plane, and to filter each triangle against $\alpha$ evaluated at its own
circumcenter (or centroid) rather than against one fixed number.

\begin{task}
Generalize your filter to accept $\alpha$ as a function $\alpha(x,y)$
instead of a constant (a simple analytic function, e.g.\ depending on
distance to a reference point, is enough --- you do not need a full mesh
size-field data structure for this homework). Test on a point cloud whose
density genuinely varies across the domain, and show that a well-chosen
$\alpha(x,y)$ recovers the shape correctly where a single constant $\alpha$
cannot.
\end{task}

\begin{note}
Export your triangulation to Gmsh's \texttt{.pos} format (or any format you
can visualize) at every stage --- after Bowyer--Watson and after alpha
filtering --- so you can \emph{see} bugs instead of guessing at them.
\end{note}

% ======================================================================
\homeworkpart{2}{Improving the Alpha Shape --- Chew's Algorithm}
% ======================================================================

Simply deleting large-circumradius triangles can leave behind a shape with a
jagged, uneven boundary: a triangle can be thin and skinny yet still have a
circumradius just under $\alpha$, so it survives the filter even though it
poorly represents the local shape. \textbf{L.~Paul Chew's algorithm}
(1989/1993) improves this by \emph{adding} well-chosen points to the
triangulation rather than only removing triangles --- and it reuses, almost
without modification, the point-insertion routine you already wrote in
Part~1.

\subsection*{2.1 --- The idea}

Call a triangle $T$ \textbf{bad} if its circumradius-to-shortest-edge ratio
exceeds a quality bound $\rho$ (a typical choice is $\rho\approx 1$, which
excludes triangles with an angle smaller than about $20^\circ$). Chew's
algorithm repeatedly picks a bad triangle and inserts a new point,
a \textbf{Steiner point}, at its \textbf{circumcenter}:

\needspace{11\baselineskip}
\begin{lstlisting}[style=code,caption={Chew's refinement algorithm},captionpos=b]
repeat:
    find a triangle T that is "bad"
        (circumradius / shortest edge > rho,  OR  circumradius > alpha)
    if none exists: stop
    c <- circumcenter(T)
    insert c into the triangulation
        (exactly the routine of Task 2: walk, cavity, retriangulate)
\end{lstlisting}

Why the circumcenter? By definition it is equidistant from all three
vertices of $T$ and lies at the centre of $T$'s (large, empty) circumcircle
--- inserting it there is the single point that most effectively breaks $T$
and its immediate neighbours into smaller, better-shaped triangles, since
after insertion no triangle incident to $c$ can have a circumcircle larger
than the one that existed before (the new circumcircles are all contained
in the old one).

\subsection*{2.2 --- Termination}

For a large enough quality bound ($\rho\gtrsim 1$, roughly a $20^\circ$
minimum-angle guarantee) Chew's algorithm is known to terminate, inserting a
number of Steiner points bounded in terms of the geometry of $S$. For this
homework you do not need to reprove that bound --- but you should still
protect your implementation with a maximum number of iterations, and you
should expect (and check) that the number of bad triangles decreases
monotonically as refinement proceeds.

\begin{note}
A circumcenter can fall \emph{outside} the current alpha shape (e.g.\ across
a concavity). A simple, adequate policy for this homework: only insert
circumcenters that fall inside the current alpha shape; a bad triangle whose
circumcenter would fall outside is instead handled by splitting its longest
boundary edge at its midpoint. You are free to justify and implement a
different policy, as long as you state it clearly.
\end{note}

\begin{task}\label{tsk:chew}
Implement Chew's refinement loop on top of your Part~1 code. After
refinement has converged, \textbf{re-apply the alpha-shape filter} of
Section~1.7: the new, smaller triangles created around every inserted
point may now individually satisfy the $\alpha$ bound even where the
original coarse triangles did not, giving a tighter, better-quality
approximation of the point cloud's shape.
\end{task}

\begin{task}
Compare, on the same test point cloud(s) as Part~1, the alpha shape before
and after Chew refinement: report the number of triangles, the minimum
angle, and the number of triangles that fail the quality bound $\rho$, in
each case.
\end{task}

% ======================================================================
\homeworkpart{3}{PFEM: A Particle Finite Element Method for Free-Surface Flows}
% ======================================================================

The \textbf{Particle Finite Element Method} (PFEM) simulates a fluid as a
cloud of Lagrangian particles (the mesh vertices) that \emph{move} with the
flow. At every time step, the fluid domain itself changes shape --- this is
exactly a free surface --- and it must be re-detected from the current
particle positions before the next physics step can be computed. This is
precisely the job of the alpha shape you built in Part~1.

\subsection*{3.1 --- The provided solver}

The Navier--Stokes finite element time step itself is not something you
need to write: it will be provided by \textbf{MigFlow}, wrapped in a small
API for this homework. From your side, this API amounts to a single
function that performs one time step on a given triangulation. Its
interface will look essentially like this:

\needspace{20\baselineskip}
\begin{lstlisting}[style=code,caption={Interface of the provided solver (prototype)},captionpos=b]
// nv, nt              : number of vertices / triangles
// coord[2*nv]          : nodal coordinates (x0,y0,x1,y1,...)
// triangles[3*nt]      : vertex indices of each triangle (flat, CCW)
// nWall, wallEdges[2*nWall]   : edges (vertex pairs) lying on rigid walls
// nFree, freeEdges[2*nFree]   : edges (vertex pairs) lying on the free surface
// vx[nv], vy[nv], p[nv]       : nodal velocities & pressure, time t (in),
//                               updated in place to time t+dt (out)
// dt                   : time step
void pfemStep(int nv, int nt,
              const double *coord, const int *triangles,
              int nWall, const int *wallEdges,
              int nFree,  const int *freeEdges,
              double *vx, double *vy, double *p,
              double dt);
\end{lstlisting}

You do not need to know --- or reimplement --- anything about the fluid
solver itself. Your job is the loop \emph{around} it.

\subsection*{3.2 --- The PFEM time-stepping loop}

\needspace{20\baselineskip}
\begin{lstlisting}[style=code,caption={PFEM driver},captionpos=b]
build the initial mesh of the fluid domain      (Part 1 + Part 2)
classify its boundary edges into wallEdges (fixed, given a priori)
                            and freeEdges (the rest of the alpha-shape boundary)
vx, vy, p <- initial conditions

for each time step:
    pfemStep(nv, nt, coord, triangles,
             nWall, wallEdges, nFree, freeEdges,
             vx, vy, p, dt)                      # provided solver

    for each vertex i:                           # move the particles
        coord[i] <- coord[i] + dt * (vx[i], vy[i])

    re-triangulate coord from scratch             # Part 1
    recompute the alpha shape                     # Part 1 + Part 2
    re-classify boundary edges into wallEdges / freeEdges

    save/plot the current state
\end{lstlisting}

\subsection*{3.3 --- Practical points}

\begin{itemize}[leftmargin=1.4em]
  \item \textbf{Walls do not move.} Points that lie on a rigid wall should
    keep their prescribed position (or slide tangentially along the wall,
    if you choose to support that) --- they must never be treated as free
    and never be deleted or displaced off the wall by remeshing.
  \item \textbf{Choosing $\alpha$.} A natural choice is a small multiple of
    the local mesh spacing; too small an $\alpha$ will erroneously cut off
    thin, physically valid necks of fluid (e.g.\ during a splash), too large
    an $\alpha$ will fail to detect a fluid separation (e.g.\ a droplet
    pinching off).
  \item \textbf{Mesh quality.} As particles move, cluster, and spread out,
    triangles degenerate. Re-running Chew's algorithm (Part~2) after every
    remeshing step keeps triangle quality under control; you may also
    choose to remove points that have become too close to a neighbour.
  \item \textbf{Time step.} You are not required to derive a rigorous CFL
    condition, but you should pick $dt$ small enough that particles move by
    a small fraction of the local mesh size in a single step, and say so in
    your report.
\end{itemize}

\begin{task}
Implement the PFEM driver loop above around the provided \lstinline{pfemStep}
function, reusing your Part~1/Part~2 code for the remeshing step.
\end{task}

\begin{task}
Demonstrate your implementation on a classic \textbf{dam-break} test case
(a rectangular column of fluid, initially at rest, collapsing under gravity
in a rigid tank): provide a short sequence of snapshots (or a short video)
showing the evolving mesh and free surface.
\end{task}

% ======================================================================
\section*{Practical Information}
% ======================================================================

\begin{itemize}[leftmargin=1.4em]
  \item \textbf{Language.} Your code must be written in \textbf{C or C++}
    (required, in particular, to link directly against the provided
    \lstinline{pfemStep} binary/object file for Part~3).
  \item \textbf{Predicates and sorting.} Using an existing, well-tested
    implementation of \lstinline{orient2d}/\lstinline{incircle} (e.g.\
    Shewchuk's robust predicates) is clearly \textbf{allowed} --- Part~1 is
    about the triangulation algorithm and its data structure, not about
    re-deriving exact arithmetic. Likewise for the Hilbert sort: you may use
    an existing implementation, e.g.\
    \url{https://git.immc.ucl.ac.be/hextreme/HXTSort}.
  \item \textbf{Deliverables.} Your source code (buildable from a clean
    checkout, with a short \texttt{README}), together with a short report
    (max.\ 6 pages) describing your data structures, any design choices left
    open in the statement, and the results asked for in Parts~1--3.
  \item \textbf{Visualization.} Choice of weapons: there are plenty of
    reasonable options, and you are free to pick whichever suits you. Both
    \textbf{ParaView} and \textbf{Gmsh} are well suited to this kind of
    work --- e.g.\ export your triangulations in Gmsh's \texttt{.pos} format
    so that you can inspect your own results at every stage of development.
  \item \textbf{Deadline.} \emph{[to be announced on Moodle]}.
  \item \textbf{Academic integrity.} This is individual work. You may
    discuss ideas with classmates, but the code and the report you submit
    must be your own.
\end{itemize}

\end{document}
